\chapter{Introduction}
\label{chap:introduction}

Typesetting systems are usually evaluated on the quality of their
output. For a collaborative editor, however, the more decisive property
is how long an author waits between pressing a key and seeing the
result. That wait is what determines whether a writing session feels
like editing a document or like submitting a job to a queue.

\section{Motivation}

A modern research document is rarely a single file written by a single
person. It is a tree of included sections, a bibliography that is
synchronised from a reference manager, a directory of generated figures,
and a version history shared between several authors
\autocite{wilson2014best}. Every one of those moving parts adds work to
the build, and the build runs far more often than anyone reads the
finished PDF.

\begin{keybox}[Why latency dominates]
In a two-hour writing session an author triggers on the order of
\num{200} rebuilds but reads the full output perhaps twice. Reducing
median build time from \qty{8}{\second} to \qty{1}{\second} gives back
roughly \qty{23}{\minute} of that session.
\end{keybox}

\section{Scope and contributions}

This report is deliberately narrow. It measures build time and nothing
else: no rendering fidelity, no editor responsiveness, no network
effects. Within that scope it makes three contributions.

\begin{enumerate}
  \item A benchmark harness that produces documents of controlled size
        and structure, described in \cref{sec:harness}.
  \item A comparison of three build strategies over eight document
        sizes, reported in \cref{chap:results}.
  \item A cost model, \cref{eq:model}, that predicts build time from
        page count and cache state within \qty{9}{\percent}.
\end{enumerate}

\section{How to read this report}

\Cref{chap:background} fixes notation and states the assumptions behind
the cost model. \Cref{chap:method} describes the harness and the
measurement protocol. \Cref{chap:results} presents the measurements, and
\cref{chap:conclusion} discusses what generalises beyond our setup.
Supporting material is in \cref{chap:appendix-a}.

\begin{remark}
Every number in this report was produced by the scripts in the
repository that accompanies it. Re-running \texttt{make bench} should
reproduce \cref{fig:latency} up to hardware differences.
\end{remark}

\section{Conventions}

Symbols are introduced where they first appear and collected in
\cref{tab:notation}. Times are wall-clock seconds measured inside the
build worker; all other quantities are dimensionless.

\begin{table}[htbp]
  \centering
  \caption{Notation used throughout the report.}
  \label{tab:notation}
  \small
  \begin{tabularx}{\textwidth}{@{}c L l@{}}
    \toprule
    \textbf{\sffamily Symbol} & \textbf{\sffamily Meaning} & \textbf{\sffamily Unit} \\
    \midrule
    $p$      & page count of the compiled document        & pages \\
    $n$      & number of source files in the project      & files \\
    $K$      & number of typesetting passes per build     & --- \\
    $\alpha$ & fixed initialisation cost of a build       & s \\
    $\beta$  & marginal typesetting cost per page         & s/page \\
    $\gamma$ & per-pass overhead independent of $p$       & s \\
    $\rho$   & fraction of pages changed since the last build & --- \\
    \bottomrule
  \end{tabularx}
\end{table}

A short note on the markup this template expects:

\begin{description}
  \item[Cross-references] Written with \texttt{cleveref}. Typing
        \verb|\cref{tab:notation}| gives \cref{tab:notation}, with the
        word \enquote{Table} supplied automatically, so renaming a
        float never leaves a stale label behind.
  \item[Citations] Written with \texttt{biblatex}. Typing
        \verb|\autocite{key}| produces \autocite{perez2007ipython}; use
        \verb|\textcite{key}| instead when the author should be part of
        the sentence, as in \textcite{hunter2007matplotlib}.
  \item[Units and numbers] Written with \texttt{siunitx}, so
        \verb|\qty{1.4}{\second}| produces \qty{1.4}{\second} and stays
        consistent across the whole document.
\end{description}
